\section{Related work}
\label{sec:related_work}

\subsection{Dataset Limitations}

Agricultural datasets typically restrict observations to fixed top-down~\cite{andrewVisualIdentificationIndividual2021a, gaoSelfSupervisionVideoIdentification2021b, xiaoCowIdentificationFreestall2022b, yuMultiCamCows2024MultiviewImage2024} or lateral~\cite{zhaoCompactLossVisual2022, liIndividualDairyCow2021} perspectives.
While some approaches focus on discriminative anatomical regions~\cite{adamSeaTurtleID2022LongspanDataset2024a, shojaeipourAutomatedMuzzleDetection2021, moskvyakLearningLandmarkGuided2020} or combine multiple static views~\cite{bergaminiMultiviewsEmbeddingCattle2018a, li2024cattle}, they lack the continuous viewpoint coverage required for geometric analysis.

In contrast, wildlife datasets~\cite{liATRWBenchmarkAmur2020, parham2017animal, atanbori2024spots, nepovinnykhSealIDSaimaaRinged2022a, dlaminiAutomatedIdentificationIndividuals2020b} offer larger viewpoint diversity but suffer from uncontrolled environments.
In these scenarios, the absence of precise annotations and the presence of occlusion or background clutter make it difficult to decouple viewpoint effects from environmental bias.

Synthetic data, while used for pose~\cite{Jiang2022BMVC, bonettoZebraPoseZebraDetection2024} or segmentation~\cite{muLearningSyntheticAnimals2020a, pengLearningPartSegmentation2024}, remains underutilized in ReID, with existing examples lacking precise viewpoint labels~\cite{picekCzechLynxDatasetIndividual2025}.

As highlighted in \Cref{tab:dataset metrics}, MOO addresses these limitations by combining controlled diversity with precise geometric metadata.

\begin{table*}[tbp]
    \centering
    \caption{
        Comparison of wildlife and cattle re-identification datasets with MOO.
        \textit{Non-Correlated Images}: independent frames reducing temporal and background bias.
        \textit{Synthetic}: includes synthetic data.
        \textit{Annotation}: viewpoint labels provided as continuous angles (Az./El.) or discrete categories.
    }
    \label{tab:dataset metrics}

    \scriptsize

    \begin{tabular}{llrrlccc}
        \toprule

        \textbf{Species}
        & \textbf{Dataset}
        & \textbf{Individuals}
        & \textbf{Images}
        & \textbf{View}
        & \textbf{Synthetic}
        & \textbf{Non-Correlated Images}
        & \textbf{Annotation} \\

        \midrule
        \multirow{6}{*}{Wildlife}

        & ATRW~\cite{liATRWBenchmarkAmur2020}
        & 107 & 3,649 & Lateral & $\circ$ & $\bullet$ & \\

        & SealID~\cite{nepovinnykhSealIDSaimaaRinged2022a}
        & 57 & 2,080 & Lateral & $\circ$ & $\bullet$ & \\

        & SeaTurtleIDHeads~\cite{adamSeaTurtleID2022LongspanDataset2024a}
        & 438 & 8,729 & Lateral & $\circ$ & $\bullet$ & \\

        & Nyala Data~\cite{dlaminiAutomatedIdentificationIndividuals2020b}
        & 237 & 1,942 & Lateral & $\circ$ & $\bullet$ & Left vs Right \\

        & LeopardID~\cite{atanbori2024spots}
        & 430 & 6,805 & Lateral & $\circ$ & $\bullet$ & \\

        & CzechLynx~\cite{picekCzechLynxDatasetIndividual2025}
        & 219+300 & 37k+100k & Lateral & $\bullet$ & $\bullet$ & \\

        \midrule
        \multirow{10}{*}{Cattle}

        & Multi-views Embedding~\cite{bergaminiMultiviewsEmbeddingCattle2018a}
        & 439 & 17,802 & Lateral & $\circ$ & $\circ$ & Front vs Side \\

        & OpenCows2020~\cite{andrewVisualIdentificationIndividual2021a}
        & 46 & 4,736 & Top & $\circ$ & $\circ$ & \\

        & Cows2021~\cite{gaoSelfSupervisionVideoIdentification2021b}
        & 182 & 13,784 & Top & $\circ$ & $\circ$ & \\

        & MuzzleCows~\cite{shojaeipourAutomatedMuzzleDetection2021}
        & 300 & 2,900 & Face & $\circ$ & $\bullet$ & \\

        & Dairy Cows 2021~\cite{liIndividualDairyCow2021}
        & 13 & 1,485 & Lateral & $\circ$ & $\bullet$ & \\

        & MVCAID100~\cite{zhaoCompactLossVisual2022}
        & 100 & 4,073 & Lateral & $\circ$ & $\circ$ & \\

        & Cows Stall Barn~\cite{xiaoCowIdentificationFreestall2022b}
        & 48 & 8,640 & Top & $\circ$ & $\circ$ & \\

        & MultiCamCows2024~\cite{yuMultiCamCows2024MultiviewImage2024}
        & 90 & 101,329 & Top & $\circ$ & $\circ$ & \\

        & CoBRA ReID~\cite{perneelDynamicMultiBehaviourOrientationInvariant2025}
        & 48 & 11,438 & Top & $\circ$ & $\circ$ & Az. Angle \\

        & \textbf{MOO}
        & \textbf{1000} & \textbf{128,000} & \textbf{All} & $\bullet$ & $\bullet$ & \textbf{Az. \& El. Angle} \\

        \bottomrule
    \end{tabular}
\end{table*}

\subsection{Viewpoint-Agnostic ReID}
\label{subsec:viewpoint-agnostic-reid}
Viewpoint change is a fundamental challenge in patterned animal ReID, as it deeply influences the perception of discriminative features.

To achieve view-invariance, approaches range from deep metric learning~\cite{zhaoCompactLossVisual2022, zhaoMultiCenterAgentLoss2022} and local descriptor alignment~\cite{chenHolsteinCattleFace2022b} to explicitly targeting discriminative body regions~\cite{shojaeipourAutomatedMuzzleDetection2021, nepovinnykhReidentificationPatternedAnimals2025c, phyoHybridRollingSkew2018, dacLivestockIdentificationUsing2022, nepovinnykhSpeciesAgnosticPatternedAnimal2024}.
Other works attempt to resolve geometric ambiguity by unwrapping 3D textures into normalized 2D maps~\cite{hiby2009tiger, algasovUnsupervisedPelagePattern2025} or employing multi-branch networks to jointly learn from multiple views~\cite{bergaminiMultiviewsEmbeddingCattle2018a, li2024cattle}.

However, unlike human ReID~\cite{sunDissectingPersonReIdentification2019a}, the animal domain lacks a fundamental quantification of how specific angular changes degrade feature perception.
We provide this foundational analysis to guide future model development.

In this work, we conduct a systematic study of the impact of viewpoint variations on patterned animals and lay the foundation for developing more robust models for real-world applications.
